\documentclass[10pt,a4paper]{amsart}
\usepackage[margin=19mm]{geometry}
\usepackage{amsmath,amssymb,mathtools}
\usepackage[scr=rsfs,cal=euler]{mathalfa}
\usepackage{xfrac}
\usepackage[unicode,hidelinks,bookmarks=false]{hyperref}
\newcommand{\ie}{i.e.\ }
\newcommand{\cf}{cf.\ }
\newcommand{\resp}{respectively\ }
\newcommand{\Subsection}{\subsection}
\providecommand{\nobreakdash}{\nobreak\hbox{-}}
\setlength{\emergencystretch}{2em}
\multlinegap0pt
\usepackage{bm}
\usepackage{bbm}
\usepackage{dsfont}
\usepackage{xspace}
\usepackage{cancel}
\usepackage{mathtools}
\usepackage{amsthm}
\makeatletter
\@ifundefined{definition}{\newtheorem{definition}{Definition}}{}
\@ifundefined{proposition}{\newtheorem{proposition}{Proposition}}{}
\makeatother
\newcommand{\D}{\mathbb{D}}
\title[Volterra type integral operator and analytic function spaces]{Volterra type integral operator and analytic function spaces}
\author{Rahim Kargar}
\date{Source: arXiv:1805.01043v3 (CC BY 4.0)}
\begin{document}
\maketitle

\noindent\emph{Source and licence.} Volterra type integral operator and analytic function spaces. Version arXiv:1805.01043v3 (CC BY 4.0), distributed under the Creative Commons Attribution 4.0 International licence, as recorded in the arXiv licence metadata for this exact version and shown on the abstract page https://arxiv.org/abs/1805.01043v3. Full rights notice: licenses/arxiv-1805.01043-CC-BY-4.0.txt.\par\smallskip

\noindent\textit{Editorial scope.} Counted unit: the Proposition whose source label is \texttt{prop:Tn\_derivatives} in the article (source0.tex lines 779--788, source label \texttt{prop:Tn\_derivatives}), delivered together with the article's own complete original proof of it (source0.tex lines 790--845, one \texttt{proof} environment of 56 lines). No mathematics is written, repaired or re-derived: statement and proof are the article's own text line for line. Required context retained uncounted: eq:Tn\_def (lines 686--713); not-tn (lines 699--701); prop:iterated-convolution (lines 740--745). Every definition, display and label of those lines is retained unchanged. Omitted from this package and not redistributed: 1--685, 702--739, 746--778, 789--789, 846--1355. Nothing inside a retained excerpt is omitted or altered: every retained body is delivered exactly as the article's lines are written, so no omission receipt of any kind accompanies this package. Citations: the retained excerpts carry no citation command at all, so no bibliography entry of the article is transcribed and no reference is redistributed. Reference adaptation: none was needed and none was made; every reference inside the delivered unit resolves inside it, so no unresolved reference or citation is produced. Printed slips kept literal: no printed word, symbol, hypothesis or number of the counted statement or of its proof was altered. Layout and class: the article's own class is \texttt{amsart} with packages including geometry, amsmath, amssymb, graphicx, tikz, pgfplots, hyperref; the wrapper is the lane's pinned \texttt{amsart} editorial wrapper at 10pt on A4, which restates the theorem-like environments the retained text needs and adds the article's own macro definitions that the retained text needs (\texttt{\textbackslash D}), with extra wrapper packages (bm, bbm, dsfont, xspace, cancel, mathtools). The delivered numbering is the unit's own. Assets: no retained span contains a figure environment, a graphics-inclusion command, a PDF-page inclusion, a table or a TikZ picture; no image file is copied into this package, the package declares no asset dependency and no image file is redistributed with it. Licence: version arXiv:1805.01043v3 of this article is distributed under the Creative Commons Attribution 4.0 International licence, as recorded in the arXiv licence metadata for this exact version and shown on the abstract page https://arxiv.org/abs/1805.01043v3. A full rights notice is delivered at licenses/arxiv-1805.01043-CC-BY-4.0.txt. Characters: every retained excerpt is pure ASCII, so the delivered engine, XeLaTeX with its default Latin Modern text font, reports no missing character.
\begin{definition}\rm
Let $f$ and $g$ be analytic in the unit disk $\mathbb{D}$ such that $fg'$ is analytic.  
Define
\begin{equation*}
T_{g,0}[f](z)=f(z)g'(z)=:h(z).
\end{equation*}
The $n$-th order Volterra-type integral operator for $n\in\mathbb{N}=\{1,2,\ldots\}$ is defined by
\begin{equation}\label{eq:Tn_def}
T_{g,n}[f](z)
:=\underbrace{\int_0^z\int_0^{t_1}\cdots\int_0^{t_{n-1}}}_{n\ \text{times}} h(t_n)\,dt_n\cdots dt_1,
\quad z\in\mathbb{D},
\end{equation}
where the integration variables satisfy
\begin{equation}\label{not-tn}
0 \le t_n \le t_{n-1} \le \dots \le t_1 \le |z| < 1.
\end{equation}
In the iterated integral \eqref{eq:Tn_def}, the notation \eqref{not-tn} is to be understood via the standard radial parametrization \(t_j=u_j z\), $j=1,2,\ldots,n$, with \(0\le u_n\le\cdots\le u_1\le1\). 

In particular, if $n=1$, then
\begin{equation*}
  T_{g,1}[f](z)=\int_0^z h(t)\,dt=\int_0^z f(t)g'(t)dt=T_g [f](z),\quad z\in \mathbb D
\end{equation*}
and if $n=2$, then
\begin{equation*}
  T_{g,2}[f](z)=\int_0^z\int_0^{t_1} h(t_2) dt_2 dt_1
  =\int_0^z (z-t)h(t)dt=\int_0^z (z-t)f(t)g'(t)dt,\quad z\in \mathbb D.
\end{equation*}
\end{definition}

\begin{proposition}\label{prop:iterated-convolution}
Let $T_{g,n}$ be defined by \eqref{eq:Tn_def} and $h=fg'$. Then the following identity holds for all $n\in\mathbb N$:
\begin{equation*}\label{eq:iterated-to-convolution}
T_{g,n}[f](z)=\frac{1}{(n-1)!}\int_{0}^{z}(z-t)^{\,n-1}h(t)\,dt,\quad z\in\mathbb{D}.
\end{equation*}
\end{proposition}

\begin{proposition}\label{prop:Tn_derivatives}
 For every $n\in\mathbb{N}$, we have
\begin{equation*}
\big(T_{g,n}[f]\big)'(z)=T_{g,n-1}[f](z),
\end{equation*}
and for every $n\in\mathbb{N}\setminus\{1\}$,
\begin{equation*}
\big(T_{g,n}[f]\big)''(z)=T_{g,n-2}[f](z).
\end{equation*}
\end{proposition}

\begin{proof}
Let \(h(t)=f(t)g'(t)\). Fix \(r\in(0,1)\) and consider \(z\) with \(|z|\le r\). For \(n\ge1\) define
\[
K(z,t):=\frac{(z-t)^{\,n-1}}{(n-1)!}\,h(t),\quad 0\le t\le z,
\]
such that by Proposition \ref{prop:iterated-convolution}
\[
T_{g,n}[f](z)=\int_0^z K(z,t)\,dt.
\]
Since \(h\) is analytic in \(\D\), it is continuous on the compact set
\(\{t: |t|\le r\}\). Hence, the function
\(K(z,t)\) is continuous on the compact set
\(\{(z,t): |z|\le r,\ 0\le t\le z\}\). Moreover, for \(n\ge2\) the partial derivative
\[
\frac{\partial}{\partial z}K(z,t)=\frac{(n-1)(z-t)^{\,n-2}}{(n-1)!}\,h(t)
=\frac{(z-t)^{\,n-2}}{(n-2)!}\,h(t)
\]
exists and is continuous on the same compact set; for \(n=1\) we have \(\partial_z K(z,t)=0\).
Therefore, by the Leibniz rule 
\[
\frac{d}{dz}\int_0^z K(z,t)\,dt
=K(z,z)+\int_0^z \frac{\partial}{\partial z}K(z,t)\,dt,
\]
and the identity is justified since the integrand and its \(z\)-derivative are continuous on the compact domain.
%Observe that
%\[
%K(z,z)=\frac{(z-z)^{\,n-1}}{(n-1)!}\,h(z)=
%\begin{cases}
%0, & n\ge2,\\[4pt]
%h(z), & n=1.
%\end{cases}
%\]

\medskip
\noindent\textbf{(i) First derivative.} For \(n\ge2\),
\[
\bigl(T_{g,n}[f]\bigr)'(z)
=0+\int_0^z \frac{(z-t)^{\,n-2}}{(n-2)!}\,h(t)\,dt
=\;T_{g,n-1}[f](z).
\]
For \(n=1\) the Leibniz rule gives
\[
\bigl(T_{g,1}[f]\bigr)'(z)=h(z)=T_{g,0}[f](z),
\]
concluding \(\bigl(T_{g,n}[f]\bigr)'=T_{g,n-1}[f]\) for all $n\in\mathbb N$.

\medskip
\noindent\textbf{(ii) Second derivative.} If \(n\ge2\) then \(n-1\ge1\), and by part (i)
\[
\bigl(T_{g,n}[f]\bigr)''(z)
=\bigl(T_{g,n-1}[f]\bigr)'(z)
= T_{g,n-2}[f](z).
\]
When \(n=2\) this reads \((T_{g,2}[f])''=T_{g,0}[f]=h\), which agrees with a direct computation.
This completes the proof.
\end{proof}

\end{document}
